% =====================================================================
%  1bat-ud6-11.tex  —  SESSIÓ 11: L'esperança matemàtica: jocs justos i loteries
%  (sense preàmbul: s'inclou des de main.tex)
% =====================================================================
\horatitol{Sessió 11 — L'esperança matemàtica: jocs justos i loteries}%
{Quan repetim moltes vegades un joc d'atzar amb premis, quant guanyem o perdem de mitjana per partida? L'esperança matemàtica hi respon.}

\subsection{Idea clau}

Tanquem la part de càlcul de la unitat amb una eina que connecta la probabilitat amb la
\textbf{presa de decisions econòmiques}: l'\textbf{esperança matemàtica}. Respon a una
pregunta molt pràctica: si un joc (o una loteria, o una aposta) es repetís moltíssimes
vegades, \emph{quant guanyaríem o perdríem de mitjana} cada vegada? És la matemàtica que
hi ha darrere de per què «la banca sempre guanya».

\begin{idea}
\textbf{L'esperança matemàtica}\\
Si un experiment té resultats amb \textbf{guany} (o pèrdua) $x_1, x_2, \dots, x_n$ i
probabilitats $p_1, p_2, \dots, p_n$, l'esperança és:
\[
E = x_1 p_1 + x_2 p_2 + \cdots + x_n p_n = \sum x_i\, p_i.
\]
\textbf{Com es llegeix:} és el guany \textbf{mitjà per partida} a la llarga.
\begin{itemize}[leftmargin=1.4em, itemsep=1pt]
  \item $E>0$ — el joc, de mitjana, ens \textbf{afavoreix}.
  \item $E<0$ — de mitjana, hi \textbf{perdem} (és el cas de loteries i casinos).
  \item $E=0$ — \textbf{joc just} (ni guanyem ni perdem a la llarga).
\end{itemize}
\textbf{Compte:} els guanys es posen amb el seu \textbf{signe} (positiu si guanyem,
negatiu si paguem o perdem).
\end{idea}

\subsection{Exemples}

\begin{exemple}[un joc amb un dau]
Paguem $\num{1}\eur$ per tirar un dau. Si surt un $6$, ens en donen $\num{6}\eur$; si
no, no ens en donen res. El \textbf{guany net} és: $+\num{5}\eur$ si surt $6$ (cobrem
$6$ però havíem pagat $1$) i $-\num{1}\eur$ si no (perdem el que vam pagar). Amb
$P(6)=\frac{1}{6}$ i $P(\text{no }6)=\frac{5}{6}$:
\[
E = (+5)\per\frac{1}{6} + (-1)\per\frac{5}{6} = \frac{5}{6}-\frac{5}{6}=0.
\]
$E=0$: és un \textbf{joc just}. A la llarga, ni guanyem ni perdem.
\end{exemple}

\begin{exemple}[per què la loteria no compensa]
Un bitllet costa $\num{1}\eur$. Hi ha un premi de $\num{1000}\eur$ amb probabilitat
$\frac{1}{\num{2000}}$ (i res en cas contrari). El guany net és $+\num{999}\eur$ si toca
($\num{1000}$ menys $\num{1}$ del bitllet) i $-\num{1}\eur$ si no:
\[
E = (+999)\per\frac{1}{\num{2000}} + (-1)\per\frac{\num{1999}}{\num{2000}}
= \frac{999-\num{1999}}{\num{2000}}=\frac{-\num{1000}}{\num{2000}}=-0,5.
\]
$E=-0,5\eur$: de mitjana, \textbf{perdem $\num{0,5}\eur$ per bitllet}. Per això la
loteria és un mal negoci a la llarga (l'organitzador es queda, de mitjana, la meitat del
preu).
\end{exemple}

\subsection{Exercicis resolts}

\begin{exresolt}[decidir si un joc convé]
En un joc paguem $\num{2}\eur$. Traiem una carta d'una baralla de $40$; si és figura
($12$ cartes), ens en donen $\num{5}\eur$; si no, res. Calcula l'esperança i digues si
convé jugar.
\solucio
\textbf{Guany net:} $+\num{3}\eur$ si surt figura ($5$ menys els $2$ pagats);
$-\num{2}\eur$ si no. Probabilitats: $P(\text{figura})=\frac{12}{40}=0,3$;
$P(\text{no})=0,7$.
\[
E = (+3)\per 0,3 + (-2)\per 0,7 = 0,9 - 1,4 = -0,5.
\]
$E=-\num{0,5}\eur$: de mitjana es perd mig euro per partida. \textbf{No convé} jugar a
la llarga.
\resposta{$E=-\num{0,5}\eur$; no convé jugar.}
\end{exresolt}

\begin{exresolt}[esperança d'una rifa solidària]
Una entitat ven $\num{500}$ números a $\num{2}\eur$ cadascun per a una rifa amb un únic
premi de $\num{300}\eur$. Quina és l'esperança per a qui compra un número? Té sentit
comprar-lo igualment?
\solucio
\textbf{Guany net:} $+\num{298}\eur$ si toca ($300$ menys els $2$); $-\num{2}\eur$ si no.
$P(\text{toca})=\frac{1}{500}$; $P(\text{no})=\frac{499}{500}$.
\[
E = 298\per\frac{1}{500} + (-2)\per\frac{499}{500}
= \frac{298-998}{500}=\frac{-700}{500}=-1,4.
\]
$E=-\num{1,4}\eur$: de mitjana es «perd» $\num{1,4}\eur$. Tot i així \textbf{pot tenir
sentit} comprar-lo: la diferència és un \textbf{donatiu} a l'entitat (la rifa és
solidària), de manera que la decisió no es pren només per l'esperança econòmica.
\resposta{$E=-\num{1,4}\eur$; econòmicament és desfavorable, però pot tenir sentit com a
donatiu solidari.}
\end{exresolt}

\subsection{Exercicis per fer}

\begin{exercicis}
\begin{exlist}
  \item En un joc, guanyes $\num{10}\eur$ amb probabilitat $0,2$ i perds $\num{3}\eur$
        amb probabilitat $0,8$. Calcula l'esperança. És favorable?
  \item Paguem $\num{1}\eur$ per tirar una moneda: si surt cara, ens en donen
        $\num{2}\eur$; si creu, res. Calcula el guany net de cada cas i l'esperança. És
        un joc just?
  \item Una loteria ven bitllets a $\num{3}\eur$ amb un premi de $\num{1500}\eur$ i
        probabilitat $\frac{1}{1000}$. Calcula l'esperança per bitllet.
  \item Un joc té tres resultats: guanyar $\num{5}\eur$ ($p=0,3$), guanyar $\num{1}\eur$
        ($p=0,2$) i perdre $\num{4}\eur$ ($p=0,5$). Calcula l'esperança.
  \item \textbf{(Interpretació)} Si un joc té $E=0$, vol dir que en una partida concreta
        no guanyaràs ni perdràs res? Què vol dir exactament «$E=0$»?
  \item \textbf{(Raonament)} Per què les empreses de loteria i els casinos dissenyen els
        seus jocs perquè l'esperança del jugador sigui \textbf{negativa}? Com es
        relaciona això amb la frase «la banca sempre guanya»?
\end{exlist}
\end{exercicis}

% ---------------------------------------------------------------------
\fullsolucions{Sessió 11}
\begin{solucions}
\begin{sollist}
  \item $E = 10\per 0,2 + (-3)\per 0,8 = 2 - 2,4 = -0,4\eur$. \textbf{No} és favorable
        (de mitjana es perden $\num{0,4}\eur$).
  \item Guany net: $+\num{1}\eur$ si cara ($2$ menys $1$), $-\num{1}\eur$ si creu.
        $E = (+1)\per 0,5 + (-1)\per 0,5 = 0$. És un \textbf{joc just}.
  \item Guany net: $+\num{1497}\eur$ si toca; $-\num{3}\eur$ si no.
        $E = 1497\per\frac{1}{1000} + (-3)\per\frac{999}{1000}=\frac{1497-2997}{1000}
        =\frac{-1500}{1000}=-1,5\eur$.
  \item $E = 5\per 0,3 + 1\per 0,2 + (-4)\per 0,5 = 1,5 + 0,2 - 2 = -0,3\eur$.
  \item No: en una partida concreta sí que pots guanyar o perdre. «$E=0$» vol dir que
        \textbf{a la llarga}, repetint moltíssimes partides, el guany mitjà per partida
        tendeix a zero (els guanys i les pèrdues es compensen).
  \item Perquè amb una esperança negativa per al jugador, \textbf{a la llarga}
        l'empresa guanya de mitjana la diferència en cada jugada; com més partides es
        juguen, més s'apropa el resultat real a aquesta mitjana favorable a la banca.
        Per això «la banca sempre guanya»: no en cada jugada, però sí en el conjunt de
        moltíssimes.
\end{sollist}
\end{solucions}
